\documentclass[11pt,a4paper]{article}

\usepackage[margin=2.5cm]{geometry}
\usepackage[T1]{fontenc}
\usepackage[utf8]{inputenc}
\usepackage{lmodern}
\usepackage{microtype}
\usepackage{amsmath,amssymb}
\usepackage{booktabs}
\usepackage{tabularx}
\usepackage{enumitem}
\usepackage{xcolor}
\usepackage[most]{tcolorbox}
\usepackage[colorlinks=true,linkcolor=blue!50!black,citecolor=blue!50!black,urlcolor=blue!50!black]{hyperref}

\newcommand{\code}[1]{\texttt{#1}}
\newcommand{\menu}[1]{\textsf{\textbf{#1}}}
\newcommand{\key}[1]{\fbox{\footnotesize\textsf{#1}}}
\newcommand{\dV}{\Delta V}
\newtcolorbox{tip}{colback=blue!4,colframe=blue!40!black,boxrule=0.4pt,arc=1pt,left=4pt,right=4pt,top=2pt,bottom=2pt,
  fonttitle=\bfseries\small,title=Tip}
\newtcolorbox{warn}{colback=orange!6,colframe=orange!60!black,boxrule=0.4pt,arc=1pt,left=4pt,right=4pt,top=2pt,
  bottom=2pt,fonttitle=\bfseries\small,title=Watch out}
\setlist{itemsep=2pt,topsep=4pt}
\renewcommand{\arraystretch}{1.15}
\sloppy

\title{Membrane Neuron Simulator\\[0.3em]\large User manual}
\author{W.\ van Vuuren}
\date{\today}

\begin{document}
\maketitle

\begin{abstract}
\noindent The Membrane Neuron Simulator computes the static equilibrium of fluid-driven soft devices: rubber
membranes between fluid chambers that push on each other, on rigid stops and on soft tubes. You draw the device
in CAD, import it as STEP, give every solid a role, and solve. This manual explains how to use the program, from
installing it to characterising a complete mechanical neuron. The formulation behind the solver (elements, pressure
loads, contact and the Newton method) is in the separate technical reference,
\code{docs/paper/membrane\_neuron\_simulator.pdf}.
\end{abstract}

\tableofcontents
\newpage

%==========================================================================================================
\section{Overview}

A mechanical neuron has two stages, and the program has one workspace (a \emph{space}) for each, plus one for
the combination, one to analyse its characterisation and one for automated studies. The spaces are the five tabs
at the top of the window:

\begin{center}
\begin{tabularx}{\textwidth}{@{}lX l@{}}
\toprule
Tab & What you do there & Saved as \\
\midrule
\menu{Inputs → pre-activation} & Input chambers push on membranes into a closed pre-activation chamber. Solve
  and sweep the inputs. & \code{.mns} \\
\menu{Pre-activation → activation} & A membrane, driven by the pressure difference $\Delta p$ across it, squeezes a
  soft tube; gas flows through the tube. Compute the \emph{activation function} once and save it as a design. &
  \code{.mad} \\
\menu{Full neuron} & Link a \code{.mns} and a \code{.mad}: the activation design replaces one neuron membrane.
  Solve and characterise the complete neuron over many parameters. & \code{.mfn} \\
\menu{Analysis} & Plot a full neuron's stored characterisation against its swept parameters (read only,
  Section~\ref{sec:analysis}). & --- \\
\menu{Agentic research study} & Let an AI design agent explore designs through the program's command line, and follow it
  live (experimental side project, Section~\ref{sec:study}). & a study folder \\
\bottomrule
\end{tabularx}
\end{center}

Each tab is a complete window with its own menus, toolbar and model, so you can have a neuron open in the first tab
and a valve in the second at the same time.

\paragraph{Units and sign conventions.} Lengths are in mm, Young's moduli in MPa, pressures in kPa and volumes in
mm$^3$. All pressures are \textbf{gauge}: the surroundings are 0~kPa. A membrane face that touches no chamber sees
0~kPa. The only exception is the flow-resistance equations of the activation space, which use SI units
(Section~\ref{sec:flowlaws}).

\paragraph{What is simulated.} Everything is quasi-static: the program finds the equilibrium shape for the given
pressures, not the motion towards it. Membranes are large-strain rubber sheets (incompressible neo-Hookean by
default); shells add bending stiffness; fluid chambers either hold a set pressure or change pressure with their
volume; rigid bodies and soft solids are kept apart by contact. The solver is a Newton--Raphson method that ramps the
pressures up in load steps.

%==========================================================================================================
\section{Installation and start-up}

The program is written in Python (3.10 to 3.12; tested with 3.11) and runs on Windows. To install it:
\begin{enumerate}
\item install Python, from Anaconda or from python.org (tick \emph{Add python.exe to PATH});
\item double-click \code{install.bat} in the \code{MembraneNeuronSimulator} folder. It creates a Python environment in
\code{\%LOCALAPPDATA\%\textbackslash MembraneNeuronSimulator\textbackslash venv} and installs the packages listed in
\code{requirements.txt} (\code{numpy}, \code{scipy}, \code{torch}, \code{gmsh}, \code{PyQt5}, \code{pyvista},
\code{pyvistaqt}, \code{matplotlib} and a few more). The first run downloads about 1\,GB; run it again to repair the
environment.
\end{enumerate}
The environment is kept in a short path outside the project folder because Windows limits paths to 260 characters.
Without the installer, \code{python -m pip install -r requirements.txt} installs the same packages into any Python
3.10--3.12; use PyQt5 only, since PySide6 or PyQt6 next to it can break the 3D view. Building PDF reports needs a
\LaTeX\ installation (MiKTeX) with \code{pdflatex} on the path, and the agentic research study needs Claude Code.

Start the program in one of these ways:
\begin{itemize}
\item double-click the desktop shortcut or \code{Membrane Neuron Simulator.bat} (it uses the environment of
\code{install.bat}, otherwise Anaconda's Python);
\item from a terminal in the \code{MembraneNeuronSimulator} folder:
\begin{verbatim}
python run_app.py                  # empty window
python run_app.py model.step       # import a STEP file
python run_app.py project.mns      # open a neuron project
python run_app.py valve.mad        # open an activation design
python run_app.py neuron.mfn       # open a full neuron
python run_app.py --analyse neuron.mfn   # ...in the Analysis tab
\end{verbatim}
\end{itemize}

The program runs without a console. If it misbehaves, errors are written to
\code{\%USERPROFILE\%\textbackslash.membrane\_neuron\_simulator\textbackslash errors.log} and \code{crash.log}.

%==========================================================================================================
\section{Preparing the CAD model}
\label{sec:cad}

The simulator does not draw geometry; it reads a STEP file in which every region is its own solid body. These rules
apply to every space.

\begin{enumerate}
\item \textbf{One solid per region.} Each membrane or shell is a thin solid; each fluid chamber is a solid that
  fills the fluid space; frames, housings and stops are solids. In Fusion~360 a multi-body part or an assembly both
  work.
\item \textbf{Bodies touch, they do not overlap.} A membrane and the chambers on either side of it share faces. A
  chamber fills exactly the space between its walls. Overlaps make the coupling detection unreliable.
\item \textbf{A membrane is a plate with one main face.} The program builds the membrane's mid-surface from its
  \emph{largest} CAD face, moved half the measured thickness into the solid. Make each membrane a single flat or
  curved face (a disc, plate or dome), not a face split into patches.
\item \textbf{A chamber must cover the membrane face it loads.} The coupling is found geometrically: points just
  beyond each membrane face are tested against every chamber solid.
\item \textbf{Name your bodies.} Part names (or Fusion~360 body names) are carried over and appear in the model
  tree. Names such as \code{Membrane\_1}, \code{Chamber\_A}, \code{Housing} or \code{TubeFluid} let the program
  guess roles for you.
\item \textbf{Export in millimetres} as STEP (\code{.step} or \code{.stp}).
\end{enumerate}

\begin{tip}
Two generated examples show a correct model: \code{python examples/make\_neuron\_step.py} writes
\code{examples/soft\_neuron.step} (a neuron), and \code{python examples/make\_activation\_step.py} writes
\code{examples/squeeze\_valve.step} (a valve with \code{TubeFluid}, \code{InletFluid} and \code{OutletFluid}).
\end{tip}

%==========================================================================================================
\section{The window}

The neuron and activation tabs share the same layout:

\begin{itemize}
\item \textbf{3D view} (centre). Left-drag rotates, right-drag or the wheel zooms, middle-drag pans. Click a part to
  select it; \key{Ctrl}+click adds parts to the selection. Clicking the same spot again selects the part
  \emph{behind} the selected one, so you can reach inner parts without hiding anything.
\item \textbf{Model tree} (left). Every solid with its role colour. Untick a part to hide it. Right-click a selection
  for \menu{Assign role}, \menu{Show only these} and \menu{Show all}. In the activation space the tree also lists
  the \menu{Flow connections}.
\item \textbf{Side tabs} (right): \menu{Part} (role and properties of the selection), extra tabs of the space,
  \menu{Solver}, \menu{Results}.
\item \textbf{Message log} (bottom). Everything the program does and finds: couplings, warnings, convergence.
  Long jobs run in the background; the \menu{Cancel} button next to the log stops them.
\item \textbf{View toolbar}: view mode, element edges, a \menu{Transparency} slider for unselected parts and a
  \menu{Section} cut (axis plus position slider).
\end{itemize}

\begin{center}
\begin{tabular}{@{}ll@{\qquad}ll@{}}
\toprule
Key & Action & Key & Action \\
\midrule
\key{Ctrl}+\key{O} & Open STEP & \key{F5} & Solve / Simulate \\
\key{Ctrl}+\key{Shift}+\key{O} & Open project / design & \key{F6} & Sweep (neuron) \\
\key{Ctrl}+\key{S} & Save & \key{1} \key{2} \key{3} & Model / Mesh / Results view \\
\key{Ctrl}+\key{Shift}+\key{S} & Save as & \key{E} & Show element edges \\
\key{Ctrl}+\key{M} & Mesh & \key{R} & Reset camera \\
\key{Ctrl}+\key{K} & Check model & \key{F1} & Quick guide \\
\bottomrule
\end{tabular}
\end{center}

%==========================================================================================================
\section{Inputs → pre-activation (the neuron)}
\label{sec:neuron}

\subsection{Workflow}

\begin{enumerate}
\item \menu{File → Open STEP} (\key{Ctrl}+\key{O}). Every solid appears in the model tree as \emph{Unassigned}.
\item \textbf{Assign roles} (Section~\ref{sec:neuronroles}). Select one or more parts and choose the role in the
  \menu{Part} tab, or right-click → \menu{Assign role}. \menu{Edit → Auto-assign roles from names} fills in
  unassigned parts from their names. To reach parts inside a housing, untick the housing, click through it, or use a
  section cut.
\item \textbf{Set the properties} of each part in the \menu{Part} tab. Thicknesses are measured from the CAD when
  the role is assigned. Selecting several parts of the same role edits them together.
\item \menu{Simulation → Check model} (\key{Ctrl}+\key{K}). This meshes the model and writes to the log which
  chamber loads which membrane, from which side, and any warnings (for example a chamber that touches no
  membrane). Fixed nodes are drawn in blue in the mesh view.
\item \menu{Simulation → Solve} (\key{F5}), then read the \menu{Results} tab (Section~\ref{sec:results}).
\item \menu{Simulation → Sweep} (\key{F6}) to map the response over one or two inputs
  (Section~\ref{sec:sweep}).
\item \menu{File → Save project} (\key{Ctrl}+\key{S}) stores roles, properties and solver settings in a
  \code{.mns} file. The STEP file is referenced, not copied; keep it next to the project.
\end{enumerate}

\subsection{Roles and properties}
\label{sec:neuronroles}

\begin{center}
\begin{tabularx}{\textwidth}{@{}l X@{}}
\toprule
Role & Simulated as \\
\midrule
Membrane & Rubber sheet on the solid's mid-surface, no bending stiffness. \\
Shell & The same, with bending stiffness (use it for stiffer, thicker sheets). \\
Rigid body & Undeformable obstacle (housing, stop). Membranes cannot pass through it. \\
Fluid chamber & Pressure on every membrane or shell face it touches. \\
Activation membrane & A membrane replaced by a pre-computed activation design (\code{.mad}); see
  Section~\ref{sec:full}. \\
Ignore / Unassigned & Left out of the simulation. \\
\bottomrule
\end{tabularx}
\end{center}

\paragraph{Membrane and shell properties.}
\begin{description}[font=\normalfont\itshape,leftmargin=1.5em]
\item[Material model] Neo-Hookean (incompressible rubber, default) or St.\ Venant--Kirchhoff (then also
  \emph{Poisson's ratio}, default 0.45).
\item[Young's modulus] Default 0.5~MPa for a membrane, 1.0~MPa for a shell. Silicone rubbers are typically
  0.05--1~MPa.
\item[Thickness] Measured from the CAD solid when the role is assigned; type a value to override it.
\item[Elements per shortest side] Mesh density along the shorter in-plane side (default 10). 8--14 is the useful
  range; keep it the same across designs you compare.
\item[Fixed edges] \emph{All boundary edges} (default) or only \emph{Edges touching rigid bodies}.
\item[Edge support] (shells) \emph{Clamped} edges keep their slope, \emph{Pinned} edges can rotate.
\item[Pre-tension] (membranes) Isotropic in-plane tension in N/mm before any pressure is applied.
\end{description}

\paragraph{Rigid body.} Only \emph{Elements per shortest side} (the surface mesh used for contact). Curved surfaces
need a finer mesh, because contact is checked at the membrane's nodes against this mesh.

\paragraph{Fluid chamber.} The \emph{Pressure model} decides how the pressure behaves; the chamber is coloured
accordingly in the 3D view.

\begin{center}
\begin{tabularx}{\textwidth}{@{}l l X@{}}
\toprule
Model & Colour & Behaviour and settings \\
\midrule
Constant pressure (input) & green & Held at \emph{Pressure}. These are the neuron's inputs. \\
Closed: ideal gas (isothermal) & blue & Sealed air, at \emph{Pressure} when sealed. Boyle's law on the gas share.
  \emph{Ghost volume} (mm$^3$) adds volume that is connected but not drawn (tubing). \emph{Incompressible fluid}
  (\%) fills part of the chamber with liquid; the colour gets darker as it rises. \\
Closed: incompressible & purple & Sealed liquid, $P = P_0 - K(V - V_\text{fluid})$. \emph{Stiffness} in kPa per \%
  volume change (default 10; water is about 22\,000; above about 1\,000 the results barely change while solving
  gets slower). \emph{Fluid volume} defaults to the body volume; less liquid than the cavity gives suction, more
  inflates it. \\
Vent (open to surroundings) & transparent & Always 0~kPa. \\
\bottomrule
\end{tabularx}
\end{center}

\subsection{Solver settings}
\label{sec:solver}

The \menu{Solver} tab holds settings that rarely need changing:
\begin{description}[font=\normalfont\itshape,leftmargin=1.5em]
\item[Load steps] (default 10) The pressures are ramped from 0 to full in this many steps. A step that fails is
  cut automatically, so this is a starting value.
\item[Max Newton iterations] (default 40) per load step.
\item[Residual tolerance] (default $10^{-8}$) relative convergence criterion.
\item[Contact stiffness] (MPa/mm) The penalty stiffness of contact. 0 (default) chooses it automatically so that the
  penetration stays at about 1--5\,\% of the thinnest part at the highest pressure.
\end{description}

\subsection{Results}
\label{sec:results}

After a solve, the view switches to \menu{Results} (\key{3}). The \menu{Results} tab offers:
\begin{itemize}
\item the contour field: displacement magnitude, displacement X, Y or Z, or area stretch (how much a membrane's area
  has grown; above about 2 the rubber is stretched far);
\item a slider through the converged load steps; step \emph{0 (start)} shows the chamber pressures on the undeformed
  shape (for example the initial suction of an under-filled liquid chamber);
\item \emph{Same colour range for all load steps}, to compare steps by eye;
\item a table of every chamber's pressure and volume change;
\item export to CSV (table) and VTK (meshes, for ParaView), and \menu{File → Screenshot}.
\end{itemize}

\subsection{Sweeps}
\label{sec:sweep}

\menu{Simulation → Sweep} (\key{F6}) opens the sweep dialog.

\begin{enumerate}
\item Choose one or two input chambers with a range (\emph{From}, \emph{To} in kPa) and a number of
  \emph{Points}. Two inputs give a grid.
\item Press \menu{Run}. Each point starts from its neighbour's solution, so a sweep is several times faster than
  separate solves.
\end{enumerate}
The table lists every chamber's pressure at every point, and the outputs of any linked activation design. The plot
shows the chosen \emph{Plot output} against the inputs (a line, or a heat map for two inputs). \menu{Export CSV…}
writes all points, with the chambers' volume changes.

\begin{tip}
To sweep more than two parameters, or other chamber values than input pressures, link the neuron in the
\menu{Full neuron} tab, run a characterisation (Section~\ref{sec:full}) and look at it in the \menu{Analysis} tab
(Section~\ref{sec:analysis}).
\end{tip}

%==========================================================================================================
\section{Pre-activation → activation (the valve)}
\label{sec:activation}

The device is a valve: the pressure difference $\Delta p$ across a membrane squeezes a soft tube, through a part
bonded to the membrane (a pusher, rigid or deformable), against a rigid body. Gas flows through the tube from a
supply to a sink. The program computes how the tube's cross-section, the mass flow and the pressures along the
tube change with $\Delta p$. You simulate it once, save it as a design (\code{.mad}) and then use it as a part in a
neuron without simulating it again.

\subsection{Building the model}

In CAD, make solids for the membrane, the pusher (touching the membrane's face and the tube), the tube, the rigid
body behind it, and the gas: one body filling the tube exactly, a supply body against the tube's inlet end and a
sink body against its outlet end. The tube needs two planar end faces normal to its axis.

\begin{center}
\begin{tabularx}{\textwidth}{@{}l X@{}}
\toprule
Role & Simulated as \\
\midrule
Membrane / Shell & As in the neuron space, clamped along its edges. Positive $\Delta p$ pushes it towards the tube.
  It is bonded to every free rigid body or solid its face touches. \\
Channel (tube) & 3D solid (quadratic 10-node tetrahedra, compressible neo-Hookean, $E$ = 0.5~MPa,
  $\nu$ = 0.45), fixed at its two end faces. \\
Fluid, Constant pressure & A supply or sink held at \emph{Pressure}. \\
Fluid, Dynamic pressure & The pressure follows from the flow. The fluid filling the tube is cut into
  \emph{Segments} (default 10) along it, each a flow resistance in series with its own pressure on the tube wall. \\
Solid & Deformable 3D part (for example a soft pusher). \emph{Support}: free (held by what it is bonded to) or
  fixed where it touches fixed rigid bodies. \\
Rigid body & \emph{Motion}: Fixed (an obstacle, such as the block the tube lies on) or Free (moves and tilts as a
  rigid body, for example a rigid pusher). \\
\bottomrule
\end{tabularx}
\end{center}

\begin{warn}
A deformable (Solid) pusher can fail to converge on round tubes. Use a free rigid pusher there.
\end{warn}

\subsection{Flow connections}

Wherever two fluid bodies touch, and wherever a dynamic fluid faces the outside, the program creates a \emph{flow
connection}. They are listed under \menu{Flow connections} in the model tree and in the \menu{Flow} tab; selecting
one highlights its contact face. Each is one of:
\begin{itemize}
\item \textbf{Opening} (no resistance) — the default between two fluids;
\item \textbf{Orifice} — a resistance given by your equation;
\item \textbf{Closed} — the default towards the outside.
\end{itemize}

\subsection{Flow resistance equations}
\label{sec:flowlaws}

A resistance equation gives the pressure drop in \textbf{Pa} for a mass flow \code{mdot} in \textbf{kg/s}, in SI
units. Available variables:

\begin{center}\small
\begin{tabular}{@{}ll@{\qquad}ll@{}}
\toprule
\code{mdot} & mass flow (kg/s) & \code{A}, \code{P} & area (m$^2$), perimeter (m) \\
\code{rho} & density at the mean pressure & \code{h}, \code{w} & height and width of the section \\
\code{rho\_up} & upstream density & \code{L} & segment length \\
\code{mu} & viscosity (Pa\,s) & \code{Dh} & hydraulic diameter $4A/P$ \\
\code{p}, \code{p\_up}, \code{p\_down} & mean, up-, downstream pressure & & \\
\bottomrule
\end{tabular}
\end{center}

A segment uses the smallest cross-section in it; a connection uses its contact face (for a typical model, the whole
bore). The defaults are
\begin{itemize}
\item segment (laminar): \code{32*mu*L*mdot/(rho*A*Dh**2)};
\item orifice: \code{(mdot/(0.61*A))**2/(2*rho\_up)}.
\end{itemize}
A smaller orifice than the contact face needs its area typed into the equation, for example
\code{(mdot/(0.61*0.05e-6))**2/(2*rho\_up)}.

\begin{tip}
Size an orifice between the resistance of the open and of the closed tube. An orifice much smaller than the
closed tube takes the whole pressure drop, and the tube then never controls the flow.
\end{tip}

\subsection{Study settings and simulating}

The \menu{Study} tab sets the $\Delta p$ sweep and the gas:
\begin{center}\small
\begin{tabular}{@{}lll@{}}
\toprule
Setting & Default & Meaning \\
\midrule
$\Delta p$ from / to, Points & 0 / 20~kPa, 21 & the pre-activation range \\
Pressure tolerance & 0.01~kPa & structure--flow iteration converged \\
Max flow iterations & 30 & per $\Delta p$ point \\
Gas constant $R$, Temperature & 287.05~J/(kg\,K), 293.15~K & ideal gas, $\rho = p_\text{abs}/(RT)$ \\
Atmospheric pressure & 101.325~kPa & gauge to absolute \\
Viscosity $\mu$ & $1.81\times10^{-5}$~Pa\,s & air \\
\bottomrule
\end{tabular}
\end{center}

Press \menu{Simulate} (\key{F5}). At every $\Delta p$ the structure and the flow are iterated until the pressures on
the tube wall stop changing; the live plot follows the newest point.

\subsection{Outputs and results}

A dynamic fluid has a list of named \emph{Outputs}: the gas pressure of one segment, for example \emph{activation}
at segment 10. Add, rename or remove them in the \menu{Part} tab; choosing a segment highlights it in yellow in the 3D
view. Changing outputs does not require a new simulation.

The \menu{Results} tab shows the minimum tube area and the mass flow against $\Delta p$, the pressures at the tube
ends and at the outputs against $\Delta p$, and the area and pressure profile along the tube at the selected point.
\menu{Save design…}, \menu{Export CSV…}, \menu{Export VTK…} and \menu{Screenshot…} are underneath.

\paragraph{Saving.} \menu{File → Save design} (\key{Ctrl}+\key{S}) stores the roles, connections, settings, the
results and the deformed shapes in one \code{.mad} file. Opening it shows the curves without simulating again. From
Python, \code{ActivationDesign.load("valve.mad")} gives \code{.area(dp)}, \code{.mass\_flow(dp)} and
\code{.end\_pressure(dp)}.

\begin{warn}
A design is valid only for the supply and sink pressures it was simulated with. Change those and simulate again.
\end{warn}

%==========================================================================================================
\section{Full neuron}
\label{sec:full}

The full neuron combines a neuron project (\code{.mns}) and an activation design (\code{.mad}). One membrane of the
neuron is replaced by the design's pre-computed curve of swept volume against $\Delta p$, so the valve is not
simulated again.

\subsection{Linking a design into a neuron}

There are two ways:
\begin{itemize}
\item \textbf{In the neuron space}: give the membrane the role \emph{Activation membrane}, choose the \code{.mad} file
  under \emph{Activation design} and, if needed, the \emph{Driving chamber} (the chamber whose pressure pushes the
  membrane towards the tube; \emph{Automatic} takes the closed chamber it touches). Solves and sweeps then also report
  the design's outputs.
\item \textbf{In the Full neuron tab} (recommended for characterisation):
  \begin{enumerate}
  \item \menu{File → Import inputs to pre-activation…} (\key{Ctrl}+\key{I}) loads a \code{.mns}; a copy is embedded
    in the full neuron, the original is not changed.
  \item \menu{File → Import pre-activation to activation…} (\key{Ctrl}+\key{Shift}+\key{I}) loads a \code{.mad}. It
    must have been simulated with a version that records swept volume; otherwise simulate it again.
  \item In the \menu{Link} tab, choose the neuron part to replace and the driving chamber.
  \end{enumerate}
\end{itemize}

Copies of both imported files are kept in the full neuron, so it keeps working if the files are moved, and later
changes to them do not reach it until you update (Section~\ref{sec:outdated}).

The model tree shows two groups, \emph{Inputs to pre-activation} and \emph{Pre-activation to activation}. Click a
group to move or rotate it in the view (display only). Only chamber values (pressures, stiffness, volumes) can be
edited here, not the chamber models or the design.

\subsection{Solving and characterising}

\menu{Solve once at the set values} (\key{F5}) solves the neuron at the current chamber values.

\menu{Run characterisation} (\key{F6}) solves a grid. In the \menu{Record \& run} tab:
\begin{enumerate}
\item add sweep axes with \menu{Add parameter}: a part, a field (any numeric chamber value), from, to and points.
  You can sweep as many as you want; the number of solves is the product of the points;
\item tick what to record: chamber pressures \code{P:<chamber>} and volume changes \code{dV:<chamber>}, the
  pre-activation \code{dp}, the design outputs \code{out:<name>}, the tube area and the mass flow;
\item optionally \emph{Store the deformed shapes}, to see every point in 3D after reopening (larger file).
\end{enumerate}
The grid is solved in a snake-like order so that each point starts from a prediction made from its neighbours.
Dense axes are cheap: most of the cost is in the first few kPa where membranes are slack.

When the run finishes, the dataset is stored in the \code{.mfn} file and the file is saved. Opening the file later
shows the stored dataset; a warning appears when the model has changed since. The \menu{Results} tab lists every
point; \menu{Export CSV…} writes the table. From Python, the stored dataset can be interpolated at any parameter
values without simulating again.

\begin{warn}
If $\Delta p$ leaves the range the activation design was simulated over, the result is \textbf{extrapolated}: rows are
marked orange, the 3D view shows a red overlay and a message lists the points. Simulate the design over a wider
$\Delta p$ range.
\end{warn}

\menu{Analyse…} in the \menu{Results} tab (\menu{Simulation → Analyse characterisation}, \key{F7}) saves the full
neuron and opens its characterisation in the \menu{Analysis} tab.

\subsection{Out-of-date components}
\label{sec:outdated}

As in a CAD assembly, the full neuron keeps using the versions of the \code{.mns} and \code{.mad} it imported. When
you save changes to one of those files (in its own tab or anywhere else), the full neuron notices within a couple of
seconds:
\begin{itemize}
\item its group in the model tree turns orange with a warning sign, and hovering it says why;
\item \menu{Update} appears in the toolbar (with the number of out-of-date components) and in the \menu{File} menu;
\item the status bar and the log name the changed file. A file that is no longer found is reported in red; the copy
  kept in the full neuron is then used.
\end{itemize}
Right-click a group for \menu{Update from <file>} or \menu{Open <file> in its tab}; clicking a group shows the same
two buttons in its properties.

Updating a \textbf{neuron} shows what changed: values changed in the \code{.mns} since it was imported are taken,
values you changed in the full neuron (fluid parameters, solver settings) are kept, and for a value changed in both
places a check box chooses which wins. If the linked membrane is no longer a membrane, the link is cleared. Updating
a \textbf{design} takes the new version at once (it must have simulated results). Results and the stored
characterisation made with the old version are then reported as outdated: run them again.

%==========================================================================================================
\section{Analysis}
\label{sec:analysis}

The \menu{Analysis} tab plots a full neuron's stored characterisation without simulating anything. Open a
\code{.mfn} with \menu{File → Open full neuron…} (\key{Ctrl}+\key{O}), or press \menu{Analyse…} in the
\menu{Full neuron} tab. The model tree and the properties show every part's settings, read only (a swept setting
shows its range and the value currently shown); change them in their own tabs.

\begin{enumerate}
\item Choose the result to plot as \emph{Z} (anything recorded, for example the pre-activation chamber's
  pressure).
\item Tick swept parameters as plot axes, in order. The first gives a curve, the second a surface, the third a row
  of surfaces, the fourth a grid of surfaces and the fifth layers of grids: move the \textbf{scroll wheel} over the
  plots to go up and down the layers.
\item From the third ticked parameter on, the parameter shows \emph{how many} of its values to use instead of a
  slider. Every unticked parameter is held at its slider's value.
\item \emph{Interpolate between the simulated points}: off, only simulated points are used and the sliders snap to
  them; on, values in between are interpolated linearly from the neighbouring points.
\end{enumerate}

All panels share one Z scale and colour bar. Red points were extrapolated beyond the activation design's range,
black crosses did not converge, and gaps were not solved. \emph{Surfaces as} switches between 3D surfaces and flat
colour maps. Surfaces are drawn by the graphics card: drag to turn them (all panels turn together), right-drag or the
scroll wheel to zoom (the wheel changes layers instead when five axes are ticked). Click a point in a plot to show the
neuron and the design deformed there in the \menu{3D view} (interpolated when interpolation is on); the point is
ringed in orange and its values are listed under \emph{Point shown in 3D}.

\begin{tip}
The program can also be started straight into this tab: \code{python run\_app.py --analyse neuron.mfn}.
\end{tip}

%==========================================================================================================
\section{Design studies with an AI agent}
\label{sec:study}

The \menu{Agentic research study} tab runs a design agent (Claude Code) on a study folder and lets you follow it. The agent
cannot write or run code: it only uses the program's \code{mns} command line, builds designs from parametric design
files, runs them, and writes a \LaTeX\ report per design and one for the study.

\begin{warn}
The agentic research study is very experimental and a side project next to the simulator itself. Expect rough edges, and check
the agent's designs and conclusions yourself before relying on them.
\end{warn}

\begin{enumerate}
\item \menu{New study…}: choose an empty folder, write the \emph{brief} (the goal, the free parameters and their
  ranges, what to measure, the time budget) and add your starting models (\code{.mns}, \code{.mad},
  \code{.mfn}). They are imported as the first designs.
\item \menu{Start agent} opens Windows Terminal in the study folder with Claude Code running. The first time, accept
  Claude Code's question whether to trust the folder.
\item Watch the tab. With \emph{Follow the agent} ticked it shows the design the agent works on, the deforming shape
  during solves, the agent's current command and its reason, the design list, the event log, and every plot and
  picture. \menu{Open in its space} opens a design in the matching tab.
\item \menu{Edit brief}, \menu{Open folder} and \menu{Study report} give quick access to the folder.
\end{enumerate}

At the end, \code{report/study.pdf} holds the overall findings. Each design has a folder
\code{designs/<ID>\_<name>/} with \code{design.yaml} (parameters, geometry, roles), \code{model.step}, the
simulator project, \code{results/}, \code{renders/} and \code{report/design.pdf}.

You can use the same command line yourself, in the study folder: \code{mns help} lists the commands, and
\code{MANUAL.md} in the study folder documents them and the design-file format. For example:
\begin{verbatim}
mns design derive N001 thin --set Membrane1.thickness=0.3 --why "thinner"
mns sweep N002 --input Input1=0:20:5 --input Input2=0:20:5
mns report study --build
\end{verbatim}

%==========================================================================================================
\section{Troubleshooting}

\begin{center}
\begin{tabularx}{\textwidth}{@{}>{\raggedright}p{0.32\textwidth} X@{}}
\toprule
Symptom & What to try \\
\midrule
A chamber ``touches no membrane'' in Check model & The chamber's faces must meet the membrane face exactly.
  Inspect with a section cut; fix overlaps or gaps in CAD. \\
A membrane's mid-surface looks wrong & The membrane must have one main face (Section~\ref{sec:cad}). \\
The solve does not reach full load & More \emph{Load steps}; fewer elements per side; check for a physical limit
  point (a membrane ballooning past its peak pressure, or snap-through). The log reports the load factor reached. \\
Very large area stretch ($>2$) & The membrane hits a wall or is over-stretched: bigger chamber, thicker or stiffer
  membrane. \\
A full neuron component is marked out of date & Its file was saved with changes: right-click the group →
  \menu{Update} (Section~\ref{sec:outdated}). \\
The tube never controls the flow & The orifice is much smaller than the closed tube; enlarge it. \\
Some $\Delta p$ points did not converge & They are marked in the plot. Raise \emph{Max flow iterations}, use a
  coarser tube mesh, or use a rigid pusher. \\
``EXTRAPOLATING'' & Simulate the activation design over a wider $\Delta p$ range. \\
The program closes or hangs & Read \code{errors.log} and \code{crash.log} in
  \code{\%USERPROFILE\%\textbackslash.membrane\_neuron\_simulator}. \\
\bottomrule
\end{tabularx}
\end{center}

\section{Known limitations}

\begin{itemize}
\item Static only: no dynamics, no time.
\item Contact is checked at mesh nodes; use a mesh fine enough to follow the obstacle's curvature. A sheet does not
  contact itself.
\item A membrane's mid-surface comes from its single largest CAD face.
\item Flow is steady, isothermal and lumped (uniform pressure per segment). No choked flow unless you write it into
  the resistance equation. A fully closed tube keeps about 2\,\% of its open area (the contact gap).
\item A deformable Solid pusher may not converge on round tubes.
\item An activation design is valid only for the supply and sink pressures it was simulated with.
\end{itemize}

\section{Files}

\begin{center}
\begin{tabularx}{\textwidth}{@{}l X@{}}
\toprule
File & Contents \\
\midrule
\code{.step} & Your CAD model (not changed by the program). \\
\code{.mns} & Neuron project: STEP path, roles, properties, solver settings, sweep settings. \\
\code{.mad} & Activation design: roles, connections, study settings, results and deformed shapes. \\
\code{.mfn} & Full neuron: embedded copies of the neuron and the design (with their file paths, to detect
  changes), the link and the stored characterisation. \\
\code{<name>.csv} & Sweep export. \\
\bottomrule
\end{tabularx}
\end{center}

\end{document}
